\documentclass[11pt]{article}
\usepackage[T1]{fontenc}
\usepackage[margin=1in]{geometry}
\usepackage{amsmath,amssymb,amsthm}
\usepackage{graphicx}
\usepackage{microtype}
\usepackage[hidelinks]{hyperref}

\newtheorem{theorem}{Theorem}
\newtheorem{lemma}{Lemma}
\newcommand{\Fp}{\mathbb F_p}

\title{A positive-characteristic counterexample to\\
finite algebraic Banach field extensions}
\author{Automated Math Discovery Run \texttt{fa\_banach\_001}\\
Model: GPT5.6}
\date{22 July 2026}

\begin{document}
\maketitle

\begin{center}
\fbox{\parbox{0.9\linewidth}{\textbf{Status.}
Counterexample likely valid, subject to human review.  This gives a full
negative answer to the exact question in Remark~7.3 of arXiv:2504.10209v1;
it does not settle the paper's broader positive-characteristic automatic-
continuity problem.}}
\end{center}

\section{Source question}

Feliks R\k{a}czka's paper \emph{Continuity of differential operators for
nonarchimedean Banach algebras} proves the following as Lemma~7.2 in
characteristic zero: if $K\subset L$ is algebraic, $K$ is a nontrivially
valued nonarchimedean field, and $L$ is a Banach $K$-algebra which is a field,
then $[L:K]<\infty$ and $L$ is nonarchimedean \cite{Raczka}.
Remark~7.3 asks whether this remains true when $\operatorname{char}K=p>0$.
It specifically observes that purely inseparable extensions defeat the
root-separation argument used in characteristic zero.

\begin{figure}[ht]
  \centering
  \includegraphics[width=0.96\linewidth]{figures/open_problem_crop.png}
  \caption{Remark~7.3 on page 11 of the source paper.}
\end{figure}

\section{Definitions and notation}

For a field $F$ of characteristic $p$, write
\[
F^p=\{x^p:x\in F\}.
\]
This is a subfield.  An extension $E/F$ is \emph{purely inseparable of
exponent one} if $x^p\in F$ for every $x\in E$.  In particular, such an
extension is algebraic, although its total degree need not be finite.

For a coefficient field $k$ and an indeterminate $t$, the Laurent-series
field $k((t))$ is equipped with the $t$-adic absolute value
\[
\left|\sum_{n\ge n_0}a_nt^n\right|=q^{\min\{n:a_n\ne0\}},
\qquad 0<q<1.
\]
It is a complete nontrivially valued nonarchimedean field.

\section{Proof intuition}

In positive characteristic, Frobenius provides exactly the phenomenon that
the characteristic-zero proof cannot see.  Start with a complete field $L$
whose coefficient field has infinitely many independent imperfection
directions, and use the Frobenius image $K=L^p$ as the base field.  The base
remains complete, while every element of $L$ becomes algebraic of degree at
most $p$ over $K$.  There is nevertheless no uniform bound on the dimension
of the whole extension: finitely many imperfection variables already produce
$p^N$ independent monomials for arbitrary $N$.

\section{Counterexample}

\begin{theorem}\label{thm:counterexample}
For every prime $p$, there are fields $K\subset L$ of characteristic $p$ such
that:
\begin{enumerate}
  \item $K$ is a complete nontrivially valued nonarchimedean field;
  \item $L$ is a field and a Banach $K$-algebra for a multiplicative
        nonarchimedean norm;
  \item $L/K$ is purely inseparable of exponent one, hence algebraic;
  \item $[L:K]=\infty$.
\end{enumerate}
Consequently, the characteristic-$p$ analogue of Lemma~7.2 in
arXiv:2504.10209 is false.  It is false even when the Banach algebra norm on
$L$ is already multiplicative and nonarchimedean.
\end{theorem}

\begin{proof}
Let $u_1,u_2,\ldots$ be algebraically independent over $\Fp$, and put
\[
k=\Fp(u_1,u_2,\ldots),\qquad L=k((t)).
\]
Give $L$ its $t$-adic absolute value, normalized by $|t|=q$ for some
$0<q<1$, and set
\[
K=L^p.
\]

We first identify $K$.  Frobenius is continuous for the $t$-adic topology,
and in characteristic $p$ it gives
\[
\left(\sum_{n\ge n_0}a_nt^n\right)^p
   =\sum_{n\ge n_0}a_n^p t^{pn}.
\]
Conversely, every coefficient in $k^p$ has a unique $p$th root in $k$.
Therefore
\begin{equation}\label{eq:frobenius-image}
K=L^p=k^p((t^p))
 =\Fp(u_1^p,u_2^p,\ldots)((t^p)).
\end{equation}
Upon writing $s=t^p$, the right-hand side is the Laurent-series field
$k^p((s))$, with $|s|=q^p$.  Thus $K$ is complete, nontrivially valued, and
nonarchimedean.  The same facts hold for $L$, and multiplicativity gives
$|ax|=|a||x|$ for $a\in K$ and $x\in L$.  Hence $L$ is a Banach
$K$-algebra.

For every $x\in L$, we have $x^p\in K$, so $x$ is a zero of
\[
X^p-x^p\in K[X].
\]
Thus $L/K$ is purely inseparable of exponent one and, in particular,
algebraic.

It remains to prove that its degree is infinite.  Fix $N\ge1$ and consider
the $p^N$ elements
\[
u_1^{e_1}\cdots u_N^{e_N},
\qquad 0\le e_i<p.
\]
We claim that they are linearly independent over $K$.  Suppose that a finite
relation
\begin{equation}\label{eq:relation}
\sum_{0\le e_1,\ldots,e_N<p}
c_{e_1,\ldots,e_N}(t^p)u_1^{e_1}\cdots u_N^{e_N}=0
\end{equation}
holds, with coefficients in $K=k^p((t^p))$.  Expand each coefficient as a
Laurent series in $t^p$.  Comparing the coefficient of every power $t^{pm}$
in \eqref{eq:relation} gives a relation over $k^p$ among the same monomials.

Those monomials are linearly independent over $k^p$.  Indeed, any proposed
finite relation uses only finitely many of the variables.  In the resulting
finite rational-function field, the monomials whose exponent in each variable
lies in $\{0,\ldots,p-1\}$ form a basis over the subfield of $p$th powers:
this follows by adjoining the variables one at a time, each satisfying
$X^p-u_i^p$ and contributing degree $p$.  Hence every Laurent-series
coefficient in \eqref{eq:relation} is zero, proving the claim.

We have shown that $[L:K]\ge p^N$ for every $N$.  Consequently
$[L:K]=\infty$, completing the counterexample.
\end{proof}

\section{What the counterexample does and does not show}

The example violates the finiteness conclusion of the proposed
characteristic-$p$ lemma.  It does not violate the other conclusion: $L$ is
itself nonarchimedean.  Thus the obstruction is not a pathological
submultiplicative Banach norm; it is the unbounded $p$-rank permitted by pure
inseparability.

This packet does not give a full positive-characteristic characterization of
automatic continuity for differential operators, nor a replacement for the
source paper's Theorems~1.2--1.3.  It settles only the precise question in
Remark~7.3.

\section{Verification and novelty status}

The proof was audited along four possible failure points: the exact Frobenius
image in \eqref{eq:frobenius-image}, completeness of $K$, algebraicity of
$L/K$, and infinite $K$-linear dimension.  Each is proved explicitly above.
No computation is used.

On 22 July 2026, a bounded novelty check searched the run's lightweight
indexes and arXiv/web results using arXiv:2504.10209, the exact paper title,
the opening phrase of Remark~7.3, the source label
\texttt{NormsOnFiniteExtensions}, the author, and the terms \emph{purely
inseparable}, \emph{$p$-degree}, and \emph{algebraic Banach field extension}.
The search found the source paper and summaries of it, but no separate work
claiming to answer Remark~7.3.  This bounded search cannot certify novelty;
human literature review is still required.

\enlargethispage{6\baselineskip}
\section{Human-review recommendation}

Review \eqref{eq:frobenius-image} and the coefficientwise independence step
in \eqref{eq:relation}.  Once those two points are accepted, all source
hypotheses and the failure of finite degree are immediate.  Recommended
verdict: \texttt{likely valid}, with a separate expert novelty check before
external dissemination.

\begin{thebibliography}{9}
\bibitem{Raczka}
F. R\k{a}czka,
\emph{Continuity of differential operators for nonarchimedean Banach
algebras}, arXiv:2504.10209v1 (2025), especially Lemma~7.2 and Remark~7.3.
\url{https://arxiv.org/abs/2504.10209}
\end{thebibliography}

\end{document}
